\begin{frame}{그림으로: 중심극한정리}
  \begin{center}
  \includegraphics[height=5.2cm]{ch03_clt_illustration.png}
  \end{center}
  \vskip0.3em
  \small
  균등분포(왼쪽, 종모양과 무관)에서 시작해도, 독립적인 값 4개·30개를
  \textbf{더할수록} 분포가 가우시안(종 모양)에 가까워진다 --- 현실의
  연속 측정값이 \textbf{수많은 작은 효과의 합}일 때 가우시안이
  자연스러운 이유.
\end{frame}
